% =====================================================================
%  Degrees-of-freedom correction: derivation (camera-ready; rebuttal to
%  Reviewer P2b3, Question 1, and the AC's Action 1).  \input from the
%  "Derivation of the Profile-Marginal Score" appendix.
% =====================================================================
\subsection{The Degrees-of-Freedom Correction}
\label{sec:supp-dof}

The score~\eqref{eq:score} uses $\nu = n_I\,d_g - n_\theta$ rather
than $n_I\,d_g$, although the partition
likelihood~\eqref{eq:joint-partition} is written for a fixed
model~$\btheta$.  We now derive this correction and state the
assumptions under which~\eqref{eq:score} is the exact log of a
well-defined score.

\paragraph{Setting.}
A minimal sample~$\cS$, $|\cS| = m$, yields the candidate
$\hat\btheta_\cS = f(\bx_\cS)$.  We assume an \emph{exact minimal
solver}, $m\,d_g = n_\theta$ (true for every model in
\Cref{tab:model-dims}), so that
\begin{equation}
	\br_i(\hat\btheta_\cS) = \bzero, \qquad i \in \cS,
	\label{eq:dof-interp}
\end{equation}
with multi-root solvers handled one root at a time.  All residuals
below are evaluated at~$\hat\btheta_\cS$.  A partition splits the
observations into inliers $\cI \supseteq \cS$ and outliers~$\cO$;
$\cR = \cI \setminus \cS$ are the $n_I - m$ non-sample inliers.
Inlier residuals follow the Gaussian and outlier residuals the uniform
component of~\eqref{eq:component-densities}.

\paragraph{The score is a predictive likelihood.}
Writing $\bx_{\setminus\cS}$ for the observations outside the sample,
the score of a candidate and partition is
\begin{equation}
	\Lambda(\cS, \cI) = p(\bx_{\setminus\cS} \mid \bx_\cS, \cS, \cI).
	\label{eq:dof-predictive}
\end{equation}
Because~\eqref{eq:dof-interp} forces the candidate to interpolate its
own sample, no non-degenerate density can be assigned to~$\bx_\cS$
under its own candidate, so conditioning on the sample is forced:
$\Lambda$ is an out-of-sample predictive score, with the status of
leave-$m$-out cross-validation, not a marginal evidence.

\paragraph{Plug-in.}
Replacing the posterior over~$\btheta$ by a point mass at the solver
output, $p(\btheta \mid \bx_\cS, \sigma^2) \approx
\delta_{\hat\btheta_\cS}(\btheta)$, while keeping a prior on the noise
scale gives
\begin{equation}
	\Lambda(\cS, \cI) \approx \int_0^\infty
	p(\bx_{\setminus\cS} \mid \hat\btheta_\cS, \sigma^2, \cI)\,
	\pi(\sigma^2)\, d\sigma^2,
	\label{eq:dof-plugin}
\end{equation}
which amounts to the working assumption
$\br_i(\hat\btheta_\cS) \mid \sigma^2 \overset{\text{iid}}{\sim}
\mathcal{N}(\bzero, \sigma^2 \mathbf{I}_{d_g})$ for $i \in \cR$.

\paragraph{Conditioning out the sample.}
Given~$\bx_\cS$, both $\hat\btheta_\cS$ and, by~\eqref{eq:dof-interp},
$\br_\cS = \bzero$ are deterministic; only~$\br_\cR$ is random.  For
every bounded~$g$,
\begin{equation}
	\mathbb{E}\bigl[g(\br_\cS, \br_\cR) \mid \bx_\cS, \sigma^2, \cI\bigr]
	= \int g(\bzero, \br_\cR)\;
	p(\br_\cR \mid \hat\btheta_\cS, \sigma^2, \cI)\, d\br_\cR .
	\label{eq:dof-condition}
\end{equation}
The sample coordinates therefore enter with weight exactly one---not
with the density value $(2\pi\sigma^2)^{-m d_g/2}$, since they were
forced to zero by the fit rather than observed there---and contribute
nothing to the residual sum of squares: $\RSS_\cI = \RSS_\cR$.

\paragraph{Likelihood.}
With $n_O = N - n_I$ outliers,
\begin{equation}
	p(\br_\cR, \br_\cO \mid \hat\btheta_\cS, \sigma^2, \cI)
	= (2\pi\sigma^2)^{-\nu/2}
	\exp\!\Bigl({-}\frac{\RSS_\cI}{2\sigma^2}\Bigr)\,
	(2a)^{-n_O d_g},
	\quad
	\nu = (n_I - m)\,d_g = n_I\,d_g - n_\theta,
	\label{eq:dof-likelihood}
\end{equation}
using $m\,d_g = n_\theta$: exactly the $n_\theta$ Gaussian
coordinates consumed by the fit are removed, while the RSS is
unchanged.

\paragraph{Marginalization.}
With $\sigma^2 \sim \IG(\alpha_0, \beta_0)$, $\alpha_0, \beta_0 > 0$,
the \invgamma integral~\eqref{eq:invgamma-integral} at
$\alpha = \alpha_0 + \nu/2$, $\beta = \beta_0 + \RSS_\cI/2$ gives
\begin{equation}
	\Lambda(\cS, \cI)
	= (2\pi)^{-\nu/2}\,(2a)^{-n_O d_g}\,
	\frac{\beta_0^{\alpha_0}}{\Gamma(\alpha_0)}\,
	\frac{\Gamma(\alpha_0 + \nu/2)}
	{(\beta_0 + \RSS_\cI/2)^{\alpha_0 + \nu/2}},
	\label{eq:dof-marginal}
\end{equation}
whose logarithm is~\eqref{eq:score} up to the prior normalizer
$\beta_0^{\alpha_0}/\Gamma(\alpha_0)$, which is common to all
partitions and candidates.  The score is finite for every partition:
$\nu = 0$ gives the all-outlier baseline $(2a)^{-(N-m)d_g}$, and
$\RSS_\cI \to 0$ stays bounded because $\beta_0 > 0$.  Letting
$\alpha_0, \beta_0 \to 0$ and dropping the improper constant, again
common to all partitions so that the ranking is preserved, recovers
the Jeffreys form~\eqref{eq:score-jeffreys},
\begin{equation}
	\Lambda_J(\cS, \cI) \;\propto\;
	\frac{\Gamma(\nu/2)}{(\pi\,\RSS_\cI)^{\nu/2}}\,(2a)^{-n_O d_g},
	\qquad \nu > 0,\ \RSS_\cI > 0.
\end{equation}

\paragraph{Scope.}
The argument requires exact minimal candidates.  When $d_g$ does not
divide~$n_\theta$ (e.g.\ a homography with one distortion parameter,
$n_\theta = 9$, $d_g = 2$), the minimal sample over-determines
the model, its residuals are not all zero, and the interpolation
step~\eqref{eq:dof-interp} fails.  After local
optimization~\eqref{eq:dof-interp} no longer holds, $\RSS_\cI =
\RSS_\cR$ fails with it, and the same count $\nu = n_I\,d_g -
n_\theta$ follows instead from the classical projection argument for
a least-squares fit on~$\cI$.
